\section{Một cái tên đọc được đáng giá bao nhiêu}
\label{sec:ch16-mot-cai-ten}

\index{sparse autoencoder!cái tên đáng giá bao nhiêu|(}%
Bài đã đo chỗ thiếu mà mục trước để lại là bài mà
mục~\ref{sec:ch15-doc-duoc-dieu-khien-duoc} đã đọc, và câu hỏi mà mục ấy giao
lại cho chương này là câu hỏi mà mục này trả lời: một khái niệm bên trong mô
hình đọc được tên, có nghĩa là gì. Chương~\ref{ch:nut-that-khai-niem} trả lời
được một nửa, rằng trong một từ điển 65\,536 nơ-ron thì 18.84 phần trăm vừa có
tên đọc được vừa ép được đầu ra đi theo cái tên ấy. Mục này đọc nửa còn lại,
tức con số ấy chắc tới đâu.

\index{sparse autoencoder!ranh giới là trung bình của chính phân bố}%
Nửa còn lại bắt đầu ở chỗ ranh giới. Bốn góc phần tư được cắt tại
\emph{trung bình của chính hai phân bố điểm số}, chứ không tại một mức tuyệt đối
nào~\autocite{p32cbmsae}. Nghĩa là \enquote{đọc được} ở đây không có nghĩa là
đọc được, mà có nghĩa là đọc được hơn mức trung bình của chính từ điển đang xét;
tương tự cho điều khiển được. Không chỗ nào trong bài đặt một mức tuyệt đối cho
tính chất nào trong hai tính chất. Điều đó không làm con số sai, nhưng nó quyết
định con số ấy trả lời được câu hỏi nào: nó so các nơ-ron của một từ điển với
nhau, không so chúng với một tiêu chuẩn.

\index{sparse autoencoder!ba mô hình, ba con số}%
Hệ quả thứ nhất của cách cắt ấy là con số đổi theo từ điển. Bài chạy cùng phép
đo trên ba mô hình phía sau, và bảng~\ref{tab:ch16-bon-goc-ba-mo-hinh} là kết
quả.

\begin{table}[htbp]
  \centering
  \caption{Bốn góc phần tư đọc được và điều khiển được, trên ba mô hình phía
  sau~\autocite{p32cbmsae}. Ranh giới của mỗi hàng là trung bình của chính hai
  phân bố điểm số trong hàng ấy, nên bốn ô của một hàng so với nhau chứ không so
  với hàng khác. Cột cuối là cột mà cả mục này đọc, và ba giá trị của nó trải từ
  11.88 tới 19.82 phần trăm, tức hàng thấp nhất kém hàng cao nhất hơn một phần
  ba. Hai hàng đầu là mô hình ngôn ngữ thị giác và hàng thứ ba là mô hình sinh
  ảnh, nên hàng thứ ba là hàng mà phép ép được đọc trong không gian nhúng ảnh
  chứ không trong không gian nhúng chữ. Mỗi hàng cộng lại đúng 100 phần trăm
  trên 65\,536 nơ-ron.}
  \label{tab:ch16-bon-goc-ba-mo-hinh}
  \small
  \begin{tabular}{@{}lrrrr@{}}
    \toprule
    & Thấp & Đọc được, & Điều khiển được, & Cao \\
    Mô hình phía sau & cả hai & không điều khiển & không đọc được & cả hai \\
    \midrule
    LLaVA, CLIP-ViT-L + Vicuna-7B   & 36.26\% & 19.87\% & 25.03\% & 18.84\% \\
    LLaVA, DINOv2-L + Gemma2-9B     & 33.07\% & 23.35\% & 23.75\% & 19.82\% \\
    UnCLIP, CLIP-ViT-L + SD 2.1     & 30.84\% & 14.53\% & 42.76\% & 11.88\% \\
    \bottomrule
  \end{tabular}
\end{table}

Bài tóm tắt ba hàng ấy bằng một câu ở phụ lục: chỉ một phần nhỏ các nơ-ron, từ
12 tới 20 phần trăm, hữu dụng cho cả hai việc~\autocite{p32cbmsae}. Câu ấy đúng
với bảng. Phần mở đầu ở trang đầu thì không in khoảng mà in một con số, khoảng
19 phần trăm, tức hàng thứ nhất~\autocite{p32cbmsae}. Sách in cả khoảng, vì hàng
thứ ba thấp hơn hàng thứ nhất tới hơn một phần ba và một con số duy nhất không
nói được điều đó.

\index{sparse autoencoder!đổi ranh giới thì đổi kết luận}%
Hệ quả thứ hai nặng hơn, và chính bài đo nó. Phụ lục đặt câu hỏi hiển nhiên: nếu
lấy một mốc khác thì sao, chẳng hạn lấy điểm điều khiển được trung bình của
chính các nơ-ron mô hình gốc, tức các nơ-ron chưa qua SAE nào.
Bài đo mốc ấy được 0.173, so với 0.232 mà hình bốn góc dùng, rồi tính lại: bốn
góc của hàng thứ nhất thành 17.86, 9.58, 43.42 và 29.14 phần
trăm~\autocite{p32cbmsae}. Con số mà cả mục~\ref{sec:ch15-doc-duoc-dieu-khien-duoc}
xoay quanh chuyển từ 18.84 lên 29.14 phần trăm, hơn một lần rưỡi, chỉ vì một
lần đổi mốc. Và mốc mới không phải mốc tùy tiện: nó hỏi một nơ-ron của từ điển
có ép đầu ra tốt hơn một nơ-ron thường hay không, thay vì hỏi nó có tốt hơn mức
trung bình của chính từ điển hay không.

\index{sparse autoencoder!hai chỗ trong bài phải sửa trước khi trích}%
Hai chỗ trong đoạn ấy sách phải sửa trước khi dùng. Chỗ thứ nhất là câu kết của
chính đoạn, và nó thuộc đúng loại mà mục~\ref{sec:ch13-tran-cao-bao-nhieu} đã
gặp, tức một ý đứng được viết trong một câu không đứng được: bài viết rằng phát
biểu ban đầu của nó, rằng SAE có tỉ lệ thấp các nơ-ron hữu dụng, \enquote{29
phần trăm}, vẫn còn đúng. Phát biểu ban đầu là phát biểu về khoảng 19 phần trăm,
ở trang đầu; 29 là con số mới. Ý thực chất thì vẫn đứng, vì 29.14 phần trăm vẫn
là một phần nhỏ, còn câu thì không, nên sách giữ phép tính đổi mốc và bỏ câu kết
ấy. Chỗ thứ hai là một chỗ vênh mà
sách không giải thích được và vì vậy ghi lại thay vì làm mượt: mốc 0.232 nói
trên là điểm điều khiển được trung bình của hình bốn góc, trong khi bảng kết quả
ở trang 7 của cùng bài cho điểm điều khiển được của toàn bộ nơ-ron SAE là 0.198
và 0.203 theo hai cách ép, và bài mô tả hình bốn góc bằng
đúng một trong hai cách ép ấy. Ba con số ấy đáng lẽ nói về cùng một đại lượng.
Bài không đối chiếu chúng ở đâu, và sách không tìm ra cách đọc nào làm chúng
khớp.

\index{sparse autoencoder!điểm điều khiển được không chuyển được}%
Hệ quả thứ ba là hệ quả mạnh nhất, và nó là một con số duy nhất. Nó không đo
trên các nơ-ron của từ điển gốc mà trên phép sửa: bài huấn luyện tầng thắt khái
niệm cho hai mô hình phía sau bằng cùng một cách, lấy 1329 khái niệm có mặt ở cả
hai bên, rồi tính hệ số tương quan Pearson giữa điểm điều khiển được của chúng ở
hai bên. Kết quả là $r = 0.06$ với giá trị $p$ bằng 0.035, và bài kết luận rằng
\tn{khả năng điều khiển}{steerability} có vẻ phụ thuộc nặng vào nhiệm vụ và mô
hình phía sau~\autocite{p32cbmsae}. Đọc con số ấy cho đúng: giá trị $p$ ở cỡ mẫu
1329 chỉ nói rằng tương quan phân biệt được với 0, tức là một phát biểu về cỡ
mẫu, còn $0.06$ là phát biểu về độ mạnh và nó gần như bằng không.

Bài tự nêu một lý do cho chỗ ấy và lý do ấy phải đi kèm con số. Hai mô hình phía
sau có sẵn không dùng cùng một bộ mã hóa CLIP, và còn dùng khác tầng lẫn khác số
token~\autocite{p32cbmsae}. Nên phép so không phải phép so hai từ điển giống hệt
nhau đặt sau hai mô hình khác nhau; hai bên khác nhau ở nhiều chỗ cùng lúc, và
$r = 0.06$ không tách được chỗ nào gây ra chỗ nào. Cái nó vẫn nói được là điều
mà mục này cần, và nói được bất kể nguyên nhân: một khái niệm điều khiển được ở
một chỗ trong ngăn xếp thì không vì thế mà điều khiển được ở chỗ khác. Sách đọc
nó ở mức ấy, và không chuyển nó thẳng sang các nơ-ron của bảng bốn góc, vì đó là
một tập khác.

\index{sparse autoencoder!điểm đọc được phụ thuộc công cụ chấm}%
Nửa \enquote{đọc được} cũng không đứng yên. Bài chạy lại điểm đọc được của cùng
những nơ-ron ấy với bảy mô hình ngôn ngữ thị giác khác nhau đặt bên trong công
cụ dò tên, và điểm số chạy từ 0.154 tới 0.220, tức đầu cao gấp khoảng 1.43 lần
đầu thấp~\autocite{p32cbmsae}. Mô hình mà bài chọn cho mọi kết quả chính lại là mô
hình cho điểm thấp nhất trong bảy. Bài kết luận rằng điều đó xác nhận việc chọn
mô hình nào không ảnh hưởng tới phép đánh giá của nó. Câu ấy đúng theo nghĩa thứ
hạng, vì phép sửa mà bài đề xuất thắng từ điển gốc dưới cả bảy công cụ chấm, và
không đúng theo nghĩa mức, vì ranh giới bốn góc là một mức chứ không phải một
thứ hạng. Sách in khoảng lệch và không mở rộng chữ của bài.

\index{sparse autoencoder!một công cụ đứng cả hai phía}%
Ba hệ quả trên đều là chuyện ranh giới và chuyện đo lại. Chỗ thứ tư thì không,
và mục~\ref{sec:ch15-doc-duoc-dieu-khien-duoc} đã nêu nó: cùng một công cụ dò
tên vừa gán tên cho nơ-ron, vừa cấp điểm đọc được, vốn là độ khớp mà nó tính ra
lúc gán, vừa cấp điểm điều khiển được, vốn là độ gần giữa đầu ra sau khi ép và
chính cái tên nó đã gán. Điều mà mục này thêm vào là bài không chỉ bỏ sót chỗ
ấy: bài chọn nó. Công trình mà bài lấy làm mốc so sánh đối chiếu đầu ra sau khi
ép với những bức ảnh làm nơ-ron sáng nhất, tức với một thứ không do công cụ dò
tên sinh ra, và bài đổi sang đối chiếu với cái tên do công cụ ấy gán, gọi phép
đo mới là vững hơn và có nền ngữ nghĩa hơn~\autocite{p32cbmsae}. Chỗ mà bài tự
nhận thì nằm ở phần giới hạn và nó là một lời trông đợi chứ không phải một phép
đo: hiệu quả của cách tiếp cận phụ thuộc vào việc công cụ dò tên gán tên có
chính xác hay không, tuy nhiên các tiến bộ tiếp theo của mô hình thị giác ngôn
ngữ có khả năng sẽ cải thiện nó~\autocite{p32cbmsae}.

\index{sparse autoencoder!cái tên đáng giá bao nhiêu|)}%
\index{khoảng trống nghiên cứu!hướng tương lai của bài không giao}%
Gộp bốn chỗ lại thì trả lời được câu mà
chương~\ref{ch:nut-that-khai-niem} giao. Một khái niệm bên trong mô hình đọc
được tên có nghĩa là một công cụ dò chưa ai kiểm đã gán cho nó một cái tên với
độ khớp cao hơn trung bình của từ điển ấy, và nghĩa ấy hết ở đó. Nó không kéo
theo rằng ép nơ-ron ấy thì đầu ra đi theo tên, vì hai tính chất tách nhau ở gần
một nửa từ điển. Nó không kéo theo rằng con số đếm được là con số ổn định, vì
đổi mốc thì nó thành 29.14 phần trăm và đổi công cụ chấm thì thang điểm giãn ra
gấp khoảng 1.43 lần. Và nó không kéo theo rằng tính chất ấy thuộc về nơ-ron, vì
$r = 0.06$ nói nó thuộc về cặp. Còn phần hướng nghiên cứu tương lai của bài thì
liệt kê những cách dựng từ điển tốt hơn, và không mục nào trong đó hỏi làm sao
biết các tên gán ra là đúng. Đó là hình dạng mà
mục~\ref{sec:ch14-dieu-chua-giai-quyet} đã gặp ba lần ở Phần~IV, gặp lại ở đây
trong chính văn liệu của lối thoát.
